	\section{Marco Teórico}
	Este capítulo reúne los fundamentos conceptuales sobre los que descansa la solución propuesta y se
	organiza en dos bloques. El primero, de la Sección~\ref{sec:arquitectura_software} a la
	Sección~\ref{sec:seguridad_software}, desarrolla los fundamentos de la plataforma entendida como
	sistema de software: los estilos arquitectónicos que justifican dividirla en servicios
	independientes, las tecnologías con que se construye y se ejecuta, y los mecanismos que protegen la
	información sensible que administra. El segundo bloque, de la
	Sección~\ref{sec:problema_clasificacion} en adelante, desarrolla los fundamentos del clasificador
	de competencia: el marco legal que define el problema, los fundamentos del Procesamiento de
	Lenguaje Natural que permiten interpretar la narrativa de las quejas, las técnicas de
	representación vectorial que traducen ese lenguaje en señales numéricas y los modelos de
	aprendizaje automático que aprenden a distinguir entre quejas competentes e incompetentes. El orden
	no es casual: el clasificador no es un sistema aislado, sino un módulo que vive dentro de la
	arquitectura descrita en el primer bloque.
	
	\subsection{Fundamentos de arquitectura de software}
	\label{sec:arquitectura_software}
	La arquitectura de software es el conjunto de decisiones estructurales de alto nivel que definen
	cómo se organiza un sistema: qué componentes lo integran, qué responsabilidad tiene cada uno y cómo
	se comunican entre sí \cite{bbva_arquitectura_software}. Esas decisiones son costosas de revertir
	una vez que el sistema está en operación, de modo que conviene tomarlas de forma explícita y
	justificada desde el diseño, y no dejarlas al azar de la implementación.
	
	\subsubsection{Atributos de calidad}
	Una arquitectura no se evalúa únicamente por las funciones que permite construir, sino por los
	atributos de calidad que habilita o dificulta \cite{utel_atributos_calidad}. Los relevantes para
	una plataforma institucional de gestión de quejas son los siguientes:
	
	\begin{itemize}
		\item \textbf{Escalabilidad:} capacidad de atender un volumen creciente de usuarios y
		peticiones agregando recursos, sin rediseñar el sistema \cite{microsoft_escalabilidad}. Se
		distingue entre escalamiento vertical ---dar más capacidad a la misma máquina--- y horizontal
		---agregar más instancias que reparten la carga--- \cite{aws_arquitectura_distribuida}.
	
		\item \textbf{Mantenibilidad:} facilidad con que el sistema puede modificarse para corregir
		defectos o incorporar funciones nuevas. Depende directamente de qué tan acopladas estén sus
		partes.
	
		\item \textbf{Disponibilidad:} proporción del tiempo en que el sistema responde correctamente.
		En un sistema distribuido, el fallo de un componente no debe propagarse al resto
		\cite{nygard_release_it}.
	
		\item \textbf{Seguridad y confidencialidad:} capacidad de proteger la información frente a
		accesos no autorizados, atributo crítico cuando los datos administrados son denuncias que
		identifican a personas.
	
		\item \textbf{Trazabilidad:} posibilidad de reconstruir qué ocurrió con cada expediente y
		cuándo, requisito tanto operativo como jurídico en el contexto de la Defensoría.
	\end{itemize}
	
	\subsubsection{Del monolito a los microservicios}
	En una arquitectura \textbf{monolítica} toda la funcionalidad se empaqueta y despliega como una
	única unidad ejecutable \cite{ibm_monolitica}. Este enfoque es simple de desarrollar y desplegar en
	etapas tempranas, pero conforme el sistema crece presenta problemas conocidos: cualquier cambio
	obliga a reconstruir y desplegar la aplicación completa, el escalamiento debe hacerse en bloque
	aunque solo un componente esté saturado, y el acoplamiento interno dificulta que varias personas
	trabajen en paralelo sin interferir entre sí.
	
	Una arquitectura de \textbf{microservicios} descompone la aplicación en un conjunto de servicios
	pequeños e independientes, cada uno con una responsabilidad acotada, su propio ciclo de despliegue
	y una interfaz bien definida para comunicarse con los demás \cite{aws_microservicios}. Cada
	servicio puede desarrollarse, desplegarse y escalarse por separado, lo que reduce el impacto de un
	cambio y permite dimensionar los recursos según la carga real de cada componente. La
	Tabla~\ref{tab:monolito_vs_microservicios} contrasta ambos estilos.
	
	\rowcolors{2}{rowblue}{white}
	
	\small
	\renewcommand{\arraystretch}{1.25}
	\setlength{\tabcolsep}{5pt}
	
	\begin{longtable}{|p{3.2cm}|p{5.4cm}|p{5.4cm}|}
		\caption{Comparación entre arquitectura monolítica y de microservicios}
		\label{tab:monolito_vs_microservicios} \\
		\hline
		\rowcolor{headerblue}
		\headercell{Criterio} & \headercell{Monolito} & \headercell{Microservicios} \\
		\hline
		\endfirsthead
	
		\hline
		\rowcolor{headerblue}
		\headercell{Criterio} & \headercell{Monolito} & \headercell{Microservicios} \\
		\hline
		\endhead
	
		Unidad de despliegue & La aplicación completa. & Cada servicio por separado. \\ \hline
		Escalamiento & En bloque, aunque solo un módulo lo requiera. & Selectivo, por servicio. \\ \hline
		Acoplamiento & Alto: los módulos comparten proceso y memoria. & Bajo: se comunican por
		interfaces explícitas. \\ \hline
		Impacto de un fallo & Puede detener toda la aplicación. & Acotado al servicio afectado si el
		diseño lo aísla. \\ \hline
		Complejidad operativa & Baja: un solo artefacto que administrar. & Alta: más artefactos, redes
		y puntos de fallo. \\ \hline
		Trabajo en equipo & Los cambios compiten sobre el mismo código. & Equipos independientes por
		servicio. \\ \hline
	\end{longtable}
	
	\normalsize
	\setlength{\tabcolsep}{6pt}
	\renewcommand{\arraystretch}{1}
	
	La elección entre ambos estilos no es absoluta. Estudios comparativos que despliegan una misma
	aplicación bajo los dos patrones muestran que los microservicios mejoran el aprovechamiento de
	recursos y la escalabilidad, pero a costa de una complejidad operativa considerablemente mayor
	\cite{villamizar_microservices}. Un sistema distribuido introduce problemas que un monolito no
	tiene ---latencia de red, fallos parciales, consistencia entre nodos--- y que deben resolverse de
	forma explícita \cite{tanenbaum_distribuidos}. Por esa razón existen posiciones intermedias, como
	mantener servicios independientes que comparten una única base de datos: una decisión que sacrifica
	autonomía de datos a cambio de simplicidad operativa, razonable a escala institucional pequeña y
	que conviene documentar como tal.
	
	\subsubsection{Arquitectura en capas}
	De forma transversal al estilo anterior, la organización interna de cada servicio suele seguir un
	esquema en capas, en el que la presentación, la lógica de negocio y el acceso a datos se separan en
	niveles con responsabilidades distintas \cite{ibm_tres_capas}. La capa de presentación expone la
	interfaz ---en un servicio web, los puntos de acceso HTTP---; la capa de negocio concentra las
	reglas del dominio; y la capa de datos se ocupa de la persistencia. La ventaja principal es que un
	cambio en una capa no obliga a modificar las demás mientras se respete el contrato entre ellas
	\cite{microsoft_ncapas}.
	
	Esta separación resulta especialmente valiosa cuando las reglas del dominio son normativas, como en
	el caso de una defensoría: al concentrarlas en una sola capa, un cambio en el reglamento se traduce
	en modificaciones localizadas y no dispersas por toda la aplicación \cite{coticolop_tres_capas}.
	Sobre esa base pueden además incorporarse políticas transversales ---calidad de servicio,
	monitoreo, control de acceso--- sin reescribir la lógica de negocio \cite{pozalujan_arquitectura}.
	
	\subsubsection{Comunicación entre servicios}
	Una vez dividido el sistema, sus partes deben comunicarse. Existen dos familias de mecanismos.
	
	La \textbf{comunicación síncrona} sigue el modelo petición-respuesta: un servicio llama a otro y
	espera el resultado antes de continuar. El estilo arquitectónico dominante para ello es REST
	(\textit{Representational State Transfer}), formulado por Fielding, que organiza la interfaz en
	torno a recursos identificados por URI y manipulados mediante los verbos del protocolo HTTP, con
	comunicación sin estado entre peticiones \cite{fielding_rest}. Su ventaja es la simplicidad
	conceptual y la interoperabilidad; su desventaja, que acopla temporalmente a los servicios, pues si
	el servicio invocado no responde, el invocador queda bloqueado o debe manejar el fallo.
	
	La \textbf{comunicación asíncrona} desacopla al emisor del receptor mediante colas o flujos de
	eventos: el servicio origen publica un evento y continúa, sin esperar a que alguien lo consuma
	\cite{stopford_event_driven}. Ofrece mayor tolerancia a fallos y mejor absorción de picos de carga,
	al precio de una complejidad notablemente mayor y de consistencia eventual en los datos.
	
	Cuando el número de servicios crece, la comunicación directa entre ellos se vuelve difícil de
	administrar y aparecen componentes de apoyo ---registro y descubrimiento de servicios,
	configuración centralizada, balanceo de carga--- que resuelven de forma estándar esa
	infraestructura transversal \cite{spring_cloud}.
	
	\subsubsection{Persistencia de datos}
	La mayoría de los sistemas de información institucionales descansa sobre bases de datos
	relacionales, que organizan la información en tablas vinculadas por claves y garantizan las
	propiedades ACID: atomicidad, consistencia, aislamiento y durabilidad \cite{ibm_postgresql}. Esas
	garantías son especialmente pertinentes cuando cada registro corresponde a un expediente con
	efectos jurídicos, donde una escritura parcial o una lectura inconsistente no son aceptables.
	
	PostgreSQL es un gestor relacional de código abierto ampliamente utilizado en producción, con
	soporte para tipos de datos binarios, transacciones concurrentes, índices avanzados y mecanismos de
	respaldo y restauración \cite{schonig_postgresql15}. El diseño de la capa de datos debe considerar,
	además, cómo evolucionará el esquema con el tiempo y cómo se mantiene la integridad cuando varios
	componentes acceden al mismo almacén, problema central en el diseño de sistemas intensivos en datos
	\cite{kleppmann_ddia}.
	
	\subsection{Tecnologías de construcción y ejecución}
	\label{sec:tecnologias}
	Definidos los principios arquitectónicos, esta sección describe las tecnologías concretas que los
	materializan, agrupadas por capa.
	
	\subsubsection{Plataforma del lado del servidor}
	Java es un lenguaje de programación orientado a objetos, compilado a un código intermedio que se
	ejecuta sobre una máquina virtual (JVM), lo que le da independencia respecto del sistema operativo
	subyacente \cite{oracle_java}. Su madurez, su tipado estático y su ecosistema de bibliotecas lo han
	consolidado en el desarrollo de sistemas empresariales de larga vida \cite{ibm_java}.
	
	Spring Boot es un marco de trabajo construido sobre el ecosistema Spring que reduce
	significativamente la configuración necesaria para levantar una aplicación: incorpora
	autoconfiguración, un servidor de aplicaciones embebido y un modelo de dependencias preempaquetadas
	que permite pasar de cero a un servicio operativo con poco código de infraestructura
	\cite{ibm_spring_boot}. Ofrece además componentes para exponer el estado interno de la aplicación
	---métricas, verificaciones de salud, información de entorno--- indispensables para operar un
	sistema distribuido sin acceso directo a un depurador \cite{alvarezcaules_actuator}.
	
	Para el componente de clasificación se emplea Python, lenguaje interpretado de propósito general
	cuya principal ventaja en este contexto es el ecosistema de bibliotecas de análisis de datos y
	aprendizaje automático que ha crecido a su alrededor \cite{aws_python}.
	
	\subsubsection{Interfaz de usuario}
	Angular es un marco de trabajo para construir aplicaciones web de página única (\textit{single-page
	applications}), en las que la navegación ocurre en el navegador sin recargar el documento completo
	\cite{angular_overview}. Está organizado en componentes reutilizables que encapsulan plantilla,
	estilos y lógica, e incorpora de forma nativa enrutamiento, gestión de formularios, cliente HTTP e
	inyección de dependencias, lo que evita ensamblar esas piezas a partir de bibliotecas sueltas
	\cite{imagina_angular}. El resultado de su compilación es un conjunto de archivos estáticos
	---HTML, CSS y JavaScript--- que cualquier servidor web puede entregar, sin necesidad de ejecutar
	código del lado del servidor para la interfaz.
	
	\subsubsection{Contenedores}
	Un contenedor empaqueta una aplicación junto con todas sus dependencias en una unidad aislada que
	se ejecuta sobre el núcleo del sistema operativo anfitrión \cite{redhat_docker}. A diferencia de
	una máquina virtual, no incluye un sistema operativo completo, por lo que arranca en segundos y
	consume una fracción de los recursos. Su aporte principal es la reproducibilidad: la misma imagen
	se ejecuta de forma idéntica en la máquina de desarrollo y en el servidor de producción, lo que
	elimina toda una clase de fallos atribuibles a diferencias de entorno.
	
	Los contenedores son, además, la unidad natural de despliegue de una arquitectura de
	microservicios: cada servicio se distribuye como una imagen propia, versionada e inmutable. Cuando
	el número de contenedores crece hasta volverse inmanejable manualmente, se recurre a plataformas de
	orquestación que automatizan su ciclo de vida, la recuperación ante fallos y el escalamiento
	\cite{burns_kubernetes}; a escalas menores, esa administración puede resolverse con guiones de
	despliegue, opción suficiente cuando el número de servicios es reducido y el equipo pequeño.
	
	\subsubsection{Proxy inverso}
	Un proxy inverso es un servidor que se coloca frente a uno o varios servicios y recibe en su nombre
	todas las peticiones del exterior, para después reenviarlas al componente que corresponda según la
	ruta solicitada \cite{nginx_org}. Concentra en un solo punto responsabilidades que de otro modo
	habría que replicar en cada servicio: terminación de conexiones cifradas, balanceo de carga,
	entrega de contenido estático y control del tráfico entrante \cite{khanal2023nginx}.
	
	Su papel es doble en una aplicación web moderna. Por un lado entrega los archivos estáticos que
	componen la interfaz; por otro, enruta las peticiones dirigidas a la API hacia el servicio
	correspondiente. Desde el punto de vista del navegador todo proviene de un mismo origen, lo que
	simplifica el modelo de seguridad del cliente, mientras que del lado del servidor la separación en
	servicios permanece intacta.
	
	\subsubsection{Control de versiones}
	Git es un sistema de control de versiones distribuido que registra el historial completo de cambios
	de un proyecto y permite que varias personas trabajen en paralelo sobre la misma base de código
	mediante ramas que después se integran \cite{github_about_git}. Su modelo distribuido implica que
	cada copia del repositorio contiene el historial íntegro, lo que aporta redundancia y permite
	trabajar sin conexión permanente al servidor central \cite{linis_git_basics}.
	
	\subsection{Seguridad aplicada al software}
	\label{sec:seguridad_software}
	Una plataforma que administra denuncias maneja, por definición, información sensible: identifica a
	quien denuncia, a quien es denunciado y describe hechos que pueden tener consecuencias
	disciplinarias o jurídicas. Los mecanismos de esta sección son, por ello, requisitos del sistema y
	no características opcionales.
	
	\subsubsection{Autenticación basada en tokens}
	En una arquitectura distribuida, mantener la sesión del usuario en la memoria del servidor obliga a
	que todas sus peticiones lleguen siempre a la misma instancia, lo que limita el escalamiento
	horizontal. La alternativa es la autenticación basada en tokens: el servicio de autenticación
	verifica las credenciales una sola vez y emite un token firmado que el cliente adjunta en cada
	petición posterior.
	
	JSON Web Token (JWT) es el formato más extendido para ello. Se compone de tres partes ---encabezado,
	carga útil y firma--- codificadas y separadas por puntos. La carga útil contiene afirmaciones sobre
	el usuario, como su identidad, su rol y la fecha de expiración del token; la firma, calculada con
	una clave secreta, permite a cualquier servicio que la conozca verificar que el contenido no fue
	alterado \cite{spring_security_jwt}. Su propiedad más útil en un sistema distribuido es que la
	validación es local: un servicio puede autorizar una petición sin consultar al servicio emisor ni
	mantener estado de sesión compartido.
	
	Esa misma propiedad tiene un costo: al ser autocontenido, un token válido no puede revocarse antes
	de su expiración sin introducir un mecanismo adicional de listas de revocación. De ahí que se
	recomiende emitirlos con tiempos de vida cortos.
	
	\subsubsection{Almacenamiento seguro de contraseñas}
	Las contraseñas nunca deben almacenarse en texto plano ni cifradas de forma reversible, pues quien
	obtenga acceso a la base de datos obtendría también las credenciales de todos los usuarios. La
	práctica correcta es guardar el resultado de una función resumen (\textit{hash}) diseñada
	específicamente para este fin, que sea deliberadamente costosa de calcular y que incorpore un valor
	aleatorio distinto por usuario (\textit{salt}), de modo que dos usuarios con la misma contraseña
	produzcan resúmenes diferentes y no puedan emplearse tablas precalculadas para revertirlos. La
	verificación en el inicio de sesión no consiste entonces en descifrar el valor almacenado, sino en
	recalcular el resumen de la contraseña recibida y compararlo con el registrado.
	
	\subsubsection{Control de acceso basado en roles}
	La autenticación resuelve quién es el usuario; la autorización, qué le está permitido hacer. El
	control de acceso basado en roles asigna a cada usuario uno o varios roles y asocia los permisos a
	esos roles y no a las personas, lo que simplifica la administración y reduce el riesgo de conceder
	privilegios por olvido. Debe aplicarse siguiendo el principio de menor privilegio: cada rol recibe
	únicamente los permisos indispensables para su función, y cada punto de acceso declara de forma
	explícita qué rol puede invocarlo. En una defensoría, esa distinción es también organizativa: quien
	recibe una queja no debe poder emitir una resolución, ni al contrario.
	
	\subsubsection{Intercambio de recursos entre orígenes}
	Los navegadores aplican la política de mismo origen, que impide a una página consultar recursos
	alojados en un dominio, protocolo o puerto distinto del suyo. El mecanismo CORS
	(\textit{Cross-Origin Resource Sharing}) permite relajar esa restricción de forma controlada
	mediante encabezados HTTP con que el servidor declara qué orígenes acepta. La recomendación es
	enumerar explícitamente los orígenes permitidos en vez de autorizar cualquiera, y tener presente
	que el uso de un proxy inverso que sirva la interfaz y la API bajo un mismo dominio hace que el
	navegador ni siquiera active esta verificación.
	
	\subsubsection{Cifrado en tránsito}
	Toda la información que viaja entre el navegador y el servidor debe protegerse con TLS
	(\textit{Transport Layer Security}), protocolo que establece un canal cifrado y autenticado sobre
	una conexión de red no confiable. Su versión 1.3 simplificó el proceso de negociación, redujo la
	latencia de establecimiento de la conexión y eliminó los algoritmos criptográficos considerados
	inseguros en versiones previas \cite{rescorla_tls13}. Sin este cifrado, las credenciales, los
	tokens de sesión y el contenido íntegro de una denuncia viajarían legibles para cualquier
	intermediario de la red.
	
	\subsection{Problema de clasificación en el contexto de la DDP}
	\label{sec:problema_clasificacion}
	El problema central de este proyecto es automatizar la decisión de admisibilidad de una queja ante
	la Defensoría de los Derechos Politécnicos, decisión que hoy depende del criterio manual del
	personal de primer contacto y genera inconsistencias operativas. Su punto de partida es el marco
	legal del Acuerdo publicado en la Gaceta Politécnica, Número Extraordinario 622
	\cite{ipn_acuerdo622}, donde se establece a la DDP como el órgano encargado de promover, proteger y
	defender los derechos de la comunidad frente a actos u omisiones de las autoridades del Instituto.
	Comprender los límites normativos de la competencia institucional es, por tanto, el primer paso para
	diseñar un clasificador que replique la lógica jurídica con precisión.
	
	El Artículo 3.\textsuperscript{o} de ese Acuerdo establece que la Defensoría es competente para
	conocer quejas relacionadas con presuntas violaciones a los derechos de la comunidad politécnica. El
	Artículo 4.\textsuperscript{o}, en cambio, delimita las materias en las que no lo es:
	
	\begin{enumerate}
		\item Afectaciones de carácter colectivo.
		\item Resoluciones disciplinarias.
		\item Derechos de naturaleza laboral.
		\item Evaluaciones académicas de profesores, comisiones dictaminadoras o consejos técnicos
		consultivos escolares, salvo que en dicha evaluación se violen notoriamente los derechos de los
		interesados.
		\item Procedimientos y resoluciones instaurados por el Órgano Interno de Control en el Instituto.
	\end{enumerate}
	
	El procedimiento de recepción descrito en el Artículo 13 obliga al personal de la Defensoría a
	determinar, en primera instancia, si la queja es competente o incompetente. Esa decisión, en
	apariencia sencilla, resulta en la práctica altamente ambigua por la forma en que están redactados
	los hechos, y su dependencia del criterio manual genera tres conflictos críticos:
	
	\begin{enumerate}
		\item El consumo excesivo de tiempo administrativo.
		\item La variabilidad de interpretaciones ante casos similares.
		\item La saturación operativa por el incremento constante de solicitudes.
	\end{enumerate}
	
	El problema se define entonces como una tarea de Procesamiento de Lenguaje Natural (PLN) y
	clasificación supervisada, en la que el sistema debe aprender los patrones lingüísticos y normativos
	que separan una queja competente de una incompetente.
	
	\subsubsection{Taxonomía de denuncias}
	Las denuncias que recibe la DDP no se clasifican solo por su contenido textual, sino por una
	estructura jerárquica de validación normativa que el sistema debe aprender a reconocer; establecer
	una taxonomía clara es, por ello, condición previa para construir cualquier modelo. Transformar una
	narrativa jurídica en un conjunto de datos procesable por algoritmos de aprendizaje automático exige
	organizar el conocimiento normativo de forma sistemática. En el contexto de la Defensoría, esa
	taxonomía descansa en dos reglas críticas:
	
	\begin{enumerate}
		\item \textbf{Pertenencia a la comunidad y jerarquía.} Define si las personas involucradas
		pertenecen a la comunidad politécnica. El sistema debe identificar al quejoso junto con su estatus
		institucional y a la autoridad a quien se imputa la violación. Sin esa validación de identidad, la
		queja carece de base para ser procesada.
	
		\item \textbf{Exclusiones del Artículo 4.\textsuperscript{o}} Constituyen el núcleo del
		clasificador. Las denuncias se categorizan según la presencia de los cinco criterios de exclusión
		señalados antes, criterios que el sistema traduce en siete campos binarios --- atributos jurídicos
		estructurados --- que operacionalizan los siguientes contrastes normativos:
	
		\begin{itemize}
			\item Naturaleza laboral frente a derechos politécnicos.
			\item Evaluación académica pura frente a vulneración de derechos dentro de la evaluación.
			\item Afectación individual frente a colectiva.
			\item Resolución disciplinaria frente a acto de autoridad.
			\item Procedimientos del Órgano Interno de Control frente a acto u omisión de autoridad
			institucional.
		\end{itemize}
	\end{enumerate}
	
	Esta estructuración logra que la clasificación no dependa únicamente del texto libre, sino de las
	reglas establecidas en el Artículo 4.\textsuperscript{o} Al término del proceso taxonómico, el
	sistema asigna a cada expediente una etiqueta binaria definitiva ---Competente o No Competente, esta
	última con canalización a otra instancia o desechamiento---, que constituye la variable de salida de
	los modelos de clasificación supervisada.
	
	\subsubsection{Clasificación binaria y lógica de admisión}
	Definida la taxonomía normativa, el siguiente paso es determinar cómo producirá el clasificador su
	veredicto. El sistema adopta un esquema de clasificación binaria, decisión fundamentada en el
	Artículo 13 del Acuerdo Número Extraordinario 622. A diferencia de los sistemas de categorización
	temática ---que discriminan entre dominios como deportes, política o finanzas---, el clasificador de
	la DDP emite un veredicto dicotómico: Competente o No Competente. Esta reducción a dos clases
	responde a la lógica de admisión procesal por tres razones:
	
	\begin{enumerate}
		\item \textbf{Naturaleza del flujo administrativo.} Dentro de la DDP, el primer contacto funciona
		como una compuerta. No existe un estado intermedio: la queja se admite para investigación o se
		desecha y canaliza a otra instancia. La clasificación binaria representa fielmente ese proceso de
		toma de decisiones.
	
		\item \textbf{Optimización del entrenamiento.} Desde la perspectiva del aprendizaje automático,
		concentrar la muestra en dos etiquetas produce mayor densidad de datos por categoría
		\cite{manning2008ir}. Dado que el volumen de expedientes históricos etiquetados es finito, un
		modelo binario resulta más robusto y menos propenso al error que uno multiclase que intentara
		identificar además el tipo específico de exclusión ---laboral, académica o disciplinaria---.
	
		\item \textbf{Reducción del falso negativo.} La lógica de admisión prioriza la protección de los
		derechos de la comunidad. Al plantear el problema como binario es posible ajustar el umbral de
		decisión ---en especial en algoritmos como SVM--- para minimizar los falsos negativos y asegurar
		que ninguna queja potencialmente competente se descarte por una clasificación errónea.
	\end{enumerate}
	
	En consecuencia, si el modelo predice la clase No Competente, el sistema asiste al personal
	sugiriendo la fundamentación legal derivada del Artículo 4.\textsuperscript{o} Si la predicción es
	Competente, el expediente fluye de forma automática hacia la etapa de radicación en la
	Subdefensoría, con lo que se cumple la meta de reducir los tiempos de respuesta institucional.
	
	\subsubsection{Ciclo de vida de los datos en el aprendizaje supervisado}
	Para que la clasificación binaria opere de manera confiable, la narrativa jurídica del quejoso debe
	recorrer un flujo técnico que la convierta en una señal numérica apta para el aprendizaje
	automático. El ciclo de vida de los datos organiza en seis etapas ese puente entre el texto libre y
	la decisión del modelo, y asegura la trazabilidad e integridad de la información jurídica en cada
	fase:
	
	\begin{enumerate}
		\item \textbf{Captura de datos.} El ciclo inicia en la interfaz de registro, donde el quejoso
		proporciona de forma estructurada su identidad institucional y, de forma no estructurada, el relato
		de los hechos. Esta etapa asegura que la entrada cumpla los requisitos mínimos de admisión del
		Artículo 13.
	
		\item \textbf{Preprocesamiento y limpieza.} La narrativa libre se somete a técnicas de PLN: se
		elimina el ruido, se normaliza el texto y se identifican los términos clave que corresponden a las
		exclusiones normativas del Artículo 4.\textsuperscript{o}
	
		\item \textbf{Vectorización y enriquecimiento.} El texto se convierte en representaciones
		vectoriales (\textit{embeddings}). En paralelo, el sistema inyecta los atributos jurídicos binarios
		extraídos del expediente y aplica el factor de amplificación $\alpha$ para dar peso a la señal
		normativa frente a la semántica.
	
		\item \textbf{Entrenamiento y ajuste.} Con un conjunto de datos histórico previamente etiquetado,
		el modelo aprende a identificar las fronteras de decisión legal. Se emplean técnicas de validación
		cruzada para medir la robustez de la generalización.
	
		\item \textbf{Inferencia y sugerencia.} Ante una queja nueva, el modelo procesa los vectores y
		genera una predicción de competencia acompañada de un nivel de probabilidad. El sistema asiste al
		personal mostrando el veredicto sugerido y la posible fundamentación legal.
	
		\item \textbf{Retroalimentación.} El ciclo se cierra cuando el personal jurídico valida o corrige
		la sugerencia del modelo. Esas decisiones se reincorporan al repositorio de datos para el
		reentrenamiento futuro, de modo que el clasificador evolucione conforme cambian los criterios
		institucionales.
	\end{enumerate}
	
	\subsection{Fundamentos de Procesamiento de Lenguaje Natural}
	El Procesamiento de Lenguaje Natural es la rama de la inteligencia artificial que estudia cómo dotar
	a las computadoras de la capacidad de leer, interpretar y generar texto humano
	\cite{jurafsky2023slp}. Su importancia radica en que el lenguaje natural ---a diferencia de los
	datos numéricos o tabulares--- es ambiguo por naturaleza, depende del contexto y está lleno de
	matices que una máquina no puede procesar de forma directa. Para que un algoritmo razone sobre un
	texto, este debe transformarse primero en una representación matemática que preserve su significado.
	
	En este proyecto el PLN cumple un rol central: las quejas que recibe la DDP llegan como narrativas
	libres, redactadas en lenguaje coloquial o semijurídico por personas sin formación legal. Ello
	introduce desafíos ausentes en dominios más estandarizados: el vocabulario especializado varía entre
	quejosos, la estructura argumentativa suele ser implícita y la clasificación jurídica correcta
	depende no solo de las palabras empleadas, sino de la relación normativa entre los hechos descritos
	y los artículos del Acuerdo 622. Dos quejas léxicamente parecidas pueden tener, por eso,
	clasificaciones opuestas según el contexto institucional que las rodea.
	
	Para atender esos desafíos el sistema requiere dos tipos de herramientas de PLN, que se desarrollan
	en las subsecciones siguientes: el fundamento matemático que permite representar documentos como
	vectores comparables y las transformaciones lingüísticas que preparan y limpian el texto antes de
	aplicar ese fundamento.
	
	\subsubsection{Álgebra lineal en espacios vectoriales de alta dimensión}
	Para que una computadora compare documentos, determine cuáles se parecen o trace una frontera que
	separe categorías, esos documentos deben existir antes como objetos matemáticos. El álgebra lineal
	provee ese lenguaje: permite representar cualquier texto como un vector ---una lista ordenada de
	números--- y medir después las relaciones entre vectores con precisión aritmética. Esta
	representación es el fundamento sobre el que operan todos los vectorizadores y clasificadores del
	sistema; sin ella no existiría mecanismo alguno para que el modelo aprenda qué quejas se parecen
	entre sí ni cómo trazar la línea que separa las competentes de las incompetentes.
	
	La idea central es que cualquier documento puede representarse como un punto en un espacio vectorial
	de alta dimensión $\mathbb{R}^{d}$, donde $d$ es el número de características que lo describen.
	Documentos semánticamente similares ---por ejemplo, dos quejas sobre acoso académico--- ocupan
	regiones cercanas de ese espacio, mientras que documentos disímiles ---una queja laboral y una queja
	por evaluación--- quedan alejados. La similitud coseno formaliza esa noción de cercanía entre dos
	documentos $\mathbf{u}$ y $\mathbf{v}$:
	
	\begin{equation}
		\operatorname{sim}(\mathbf{u}, \mathbf{v}) =
		\frac{\mathbf{u} \cdot \mathbf{v}}{\lVert \mathbf{u} \rVert \, \lVert \mathbf{v} \rVert} =
		\frac{\displaystyle\sum_{i=1}^{d} u_i v_i}
		     {\sqrt{\displaystyle\sum_{i=1}^{d} u_i^{2}} \; \sqrt{\displaystyle\sum_{i=1}^{d} v_i^{2}}}
		\label{eq:similitud_coseno}
	\end{equation}
	
	donde $\operatorname{sim} \in [-1, 1]$. Un valor cercano a $1$ indica alta similitud semántica; esta
	métrica es la base del razonamiento por proximidad que subyace a todos los vectorizadores del
	sistema.
	
	El número de dimensiones con que se representan los documentos no es arbitrario. En este proyecto se
	trabaja con dos:
	
	\begin{itemize}
		\item \textbf{300 dimensiones}, correspondientes a spaCy. Sus vectores fueron entrenados por
		Explosion AI con word2vec, algoritmo que Mikolov et al. publicaron originalmente con vectores de
		300 dimensiones \cite{mikolov2013word2vec}, medida que terminó por convertirse en un estándar de
		la industria.
	
		\item \textbf{768 dimensiones}, correspondientes a Sentence-BERT. Ese número proviene de la
		arquitectura BERT base, cuyas capas de atención emplean una dimensión oculta de 768; al construir
		SBERT sobre BERT, el vector de salida del \textit{pooling} final hereda esa dimensión
		\cite{reimers2019sbert}.
	\end{itemize}
	
	A mayor dimensión, el modelo captura matices semánticos más finos ---por ejemplo, distinguir entre
	``no se me permitió acceder'' y ``accedí sin autorización''---, pero también exige más capacidad
	computacional y más datos para evitar la maldición de la dimensionalidad: conforme crece $d$, el
	volumen del espacio aumenta de forma exponencial y los datos se vuelven dispersos, lo que dificulta
	encontrar vecinos significativos. Sobre estos espacios el sistema ejecuta cuatro operaciones
	elementales del álgebra lineal:
	
	\begin{itemize}
		\item \textbf{Suma de vectores:} combina las representaciones de tokens individuales en una
		representación de documento.
		\item \textbf{Promedio ponderado:} base del vectorizador de spaCy.
		\item \textbf{Producto punto y norma:} los emplea internamente la SVM para calcular márgenes de
		separación.
		\item \textbf{Concatenación horizontal:} operación central de la arquitectura híbrida, que une los
		\textit{embeddings} de texto con los atributos jurídicos.
	\end{itemize}
	
	Para que estas operaciones produzcan vectores con significado real, el texto de entrada debe estar
	limpio y estandarizado. Ningún modelo puede extraer un vector confiable de una queja plagada de
	errores tipográficos, abreviaturas inconsistentes o sinónimos tratados como palabras distintas. De
	ahí la necesidad del preprocesamiento que se describe a continuación.
	
	\subsubsection{Preprocesamiento de la queja}
	Antes de que el álgebra lineal opere sobre el texto, la narrativa libre del quejoso pasa por una
	cadena de transformaciones que reducen el ruido y estandarizan la representación lingüística. Este
	preprocesamiento resulta especialmente relevante en textos jurídicos, donde las variaciones
	ortográficas, los errores tipográficos o los sinónimos coloquiales de términos normativos afectan
	directamente la calidad de los \textit{embeddings}. Las transformaciones se aplican en el siguiente
	orden.
	
	\paragraph{Tokenización.} Segmenta el texto continuo en unidades mínimas de significado, o tokens
	\cite{jurafsky2023slp}. spaCy realiza una tokenización sensible al contexto que distingue, por
	ejemplo, contracciones y signos de puntuación:
	
	\begin{center}
		\texttt{"me negó el acceso"} $\longrightarrow$
		[\,\texttt{me}, \texttt{negó}, \texttt{el}, \texttt{acceso}\,]
	\end{center}
	
	\paragraph{Lematización.} Reduce cada token a su forma canónica, o lema, y elimina las variaciones
	morfológicas \cite{jurafsky2023slp}:
	
	\begin{center}
		\texttt{negaron} $\rightarrow$ \texttt{negar}, \quad
		\texttt{discriminado} $\rightarrow$ \texttt{discriminar}, \quad
		\texttt{estudiantiles} $\rightarrow$ \texttt{estudiantil}
	\end{center}
	
	\paragraph{Filtrado de palabras vacías.} Las palabras vacías (\textit{stop words}) son términos de
	alta frecuencia y bajo contenido semántico: artículos, preposiciones y conjunciones. En textos
	jurídicos, sin embargo, algunas palabras que en otros contextos serían prescindibles ---como
	``no'', ``sin'' o ``nunca''--- pueden resultar jurídicamente decisivas. spaCy permite personalizar
	esa lista para el dominio específico de la Defensoría.
	
	\paragraph{Análisis morfosintáctico con spaCy.} El modelo \texttt{es\_core\_news\_lg} realiza,
	además de la tokenización y la lematización, un análisis morfosintáctico completo que asigna a cada
	token su categoría gramatical (\textit{POS tagging}) y sus dependencias sintácticas
	\cite{honnibal2020spacy}. Ello permite identificar, por ejemplo, que en la oración ``El director me
	negó el acceso al laboratorio'' el sujeto es \textit{director} (la autoridad), el verbo es
	\textit{negar} (el acto u omisión) y el complemento es \textit{acceso al laboratorio} (el derecho
	afectado). Esta información estructural enriquece la representación semántica más allá del léxico
	superficial y sienta las bases de la etapa de vectorización.
	
	\subsection{Representación vectorial del texto}
	Normalizado el texto por el preprocesamiento, el siguiente desafío es convertirlo en
	representaciones numéricas que los algoritmos de aprendizaje automático puedan procesar. La
	representación vectorial es, por tanto, el puente entre el lenguaje natural y los clasificadores. A
	continuación se presenta la evolución de estas técnicas, desde los modelos estadísticos hasta los
	modelos de lenguaje contextuales, mostrando en cada etapa por qué los enfoques más sencillos
	resultan insuficientes para el dominio jurídico de la DDP.
	
	\subsubsection{Modelos de bolsa de palabras y TF-IDF}
	Los modelos estadísticos de representación textual son el punto de partida histórico de la
	vectorización y permiten comprender, por contraste, las limitaciones que motivan el uso de
	\textit{embeddings} semánticos. Los modelos de bolsa de palabras (\textit{Bag of Words}, BoW)
	representan un documento como un vector de frecuencias de términos sobre el vocabulario total del
	corpus: cada palabra se convierte en una dimensión separada del espacio vectorial
	\cite{manning2008ir}, de modo que un conjunto de texto con $n$ palabras produce un espacio de $n$
	dimensiones, y la posición de un documento en él queda determinada por cuántas veces aparece cada
	término. Formalmente:
	
	\begin{equation}
		\mathbf{v}_{\mathrm{BoW}}(d) = \bigl[\, f(t_1, d),\; f(t_2, d),\; \dots,\; f(t_V, d) \,\bigr]
		\label{eq:bow}
	\end{equation}
	
	donde $f(t_i, d)$ es la frecuencia del término $t_i$ en el documento $d$ y $V$ es el tamaño del
	vocabulario. TF-IDF (\textit{Term Frequency\,--\,Inverse Document Frequency}) refina este enfoque al
	ponderar cada término según su frecuencia en el documento y su rareza en el corpus
	\cite{manning2008ir}:
	
	\begin{equation}
		\operatorname{tf-idf}(t, d, D) =
		\underbrace{\frac{f(t, d)}{\sum_{t'} f(t', d)}}_{\text{TF}} \times
		\underbrace{\log \frac{\lvert D \rvert}{\lvert \{ d' \in D : t \in d' \} \rvert}}_{\text{IDF}}
		\label{eq:tfidf}
	\end{equation}
	
	\paragraph{Limitaciones críticas para el dominio jurídico.} Aunque son eficientes en términos
	computacionales, estos modelos resultan inadecuados para clasificar las quejas de la DDP por dos
	razones:
	
	\begin{enumerate}
		\item \textbf{Pérdida del orden y del contexto.} Las frases ``el director negó el acceso'' y ``el
		acceso negó al director'' generan el mismo vector pese a significar cosas radicalmente distintas.
	
		\item \textbf{Sesgo léxico superficial.} El término ``calificación'' presenta alta frecuencia tanto
		en quejas por evaluación académica (No Competente) como en quejas por extorsión académica
		(Competente), lo que vuelve imposible distinguirlas correctamente.
	\end{enumerate}
	
	Estos modelos se exponen aquí como antecedente necesario para contextualizar la superioridad de los
	enfoques basados en \textit{embeddings} que emplea el sistema.
	
	\subsubsection{\textit{Word embeddings} estáticos: spaCy (es\_core\_news\_lg)}
	Dado que los modelos estadísticos pierden el contexto semántico, la siguiente generación de técnicas
	representa cada palabra como un vector denso aprendido de grandes corpus. Esta representación,
	conocida como \textit{embedding} estático, resuelve la limitación semántica de BoW al situar
	palabras relacionadas ---como \textit{acoso} y \textit{hostigamiento}--- en regiones cercanas del
	espacio vectorial. Los \textit{embeddings} estáticos representan cada palabra como un vector denso
	en $\mathbb{R}^{300}$, entrenado sobre grandes corpus en español mediante el algoritmo word2vec
	\cite{mikolov2013word2vec}.
	
	El modelo \texttt{es\_core\_news\_lg} de spaCy \cite{honnibal2020spacy} fue entrenado sobre el
	corpus \textit{Spanish Web Crawl} y genera vectores de 300 dimensiones por token. La representación
	de un documento completo se obtiene promediando los vectores de los tokens informativos:
	
	\begin{equation}
		\mathbf{v}_{\mathrm{spaCy}}(d) = \frac{1}{\lvert T \rvert} \sum_{t \in T} \mathbf{w}_t,
		\qquad
		T = \bigl\{\, t \in d \;\big|\; \texttt{has\_vector}(t) \wedge \neg\, \texttt{is\_stop}(t) \,\bigr\}
		\label{eq:spacy_promedio}
	\end{equation}
	
	donde $\mathbf{w}_t \in \mathbb{R}^{300}$ es el vector del token $t$ y $T$ es el conjunto de tokens
	informativos, es decir, aquellos que cuentan con vector disponible y no son palabras vacías.
	
	Al promediar todos los vectores del documento, no obstante, se pierden el orden sintáctico y el
	contexto local. Las frases ``el profesor me acosó condicionando la calificación'' y ``el profesor
	condicionó la calificación sin acosar'' generarían vectores cercanos pese a tener clasificaciones
	jurídicas distintas. Esta limitación motiva el uso de modelos contextuales como Sentence-BERT.
	
	\subsubsection{\textit{Sentence transformers}: SBERT paraphrase-multilingual-mpnet-base-v2}
	Para superar la insensibilidad al orden y al contexto de los \textit{embeddings} estáticos se
	requiere un modelo que genere una representación de toda la oración considerando las relaciones
	entre sus partes. Sentence-BERT (SBERT) satisface ese requisito y representa el estado del arte en
	codificación semántica de textos \cite{reimers2019sbert}: a diferencia de los \textit{embeddings}
	estáticos, genera un único vector para la oración completa y captura las relaciones entre palabras a
	lo largo de todo el texto mediante la arquitectura \textit{Transformer}.
	
	\paragraph{Arquitectura Transformer.} El bloque central del \textit{Transformer} es el mecanismo de
	atención multicabeza (\textit{multi-head self-attention}), que calcula para cada token su relevancia
	en función de todos los demás tokens del documento:
	
	\begin{equation}
		\operatorname{Attention}(Q, K, V) =
		\operatorname{softmax}\!\left( \frac{Q K^{\top}}{\sqrt{d_k}} \right) V
		\label{eq:atencion}
	\end{equation}
	
	donde $Q$, $K$ y $V$ son las matrices de consulta, clave y valor derivadas de los
	\textit{embeddings} de entrada, y $d_k$ es la dimensión de la clave \cite{vaswani2017attention}.
	Este mecanismo permite que el modelo asigne mayor peso a los tokens contextualmente relevantes. Por
	ejemplo, en la queja ``me reprobó porque no acepté salir con él'', la atención aprende a conectar
	\textit{reprobó} con \textit{acepté salir} e identifica que la relación causal entre ambos
	constituye un acto de acoso y no una evaluación académica pura.
	
	\paragraph{Entrenamiento contrastivo de SBERT.} SBERT extiende BERT con una capa de \textit{pooling}
	sobre los estados ocultos y se entrena mediante aprendizaje contrastivo con pares de oraciones
	semánticamente equivalentes y no equivalentes:
	
	\begin{equation}
		\mathcal{L}_{\mathrm{contrastiva}} =
		\sum_{(i,j)} \bigl( \operatorname{sim}(\mathbf{e}_i, \mathbf{e}_j) - y_{ij} \bigr)^{2}
		\label{eq:contrastiva}
	\end{equation}
	
	donde $y_{ij} = 1$ si las oraciones $i$ y $j$ son semánticamente equivalentes y $y_{ij} = 0$ en caso
	contrario.
	
	\paragraph{Modelo empleado.} El modelo \texttt{paraphrase-multilingual-mpnet-base-v2} fue entrenado
	sobre más de 50 idiomas, entre ellos el español. Genera \textit{embeddings} de 768 dimensiones por
	documento \cite{hf_mpnet_multilingual} y ha mostrado un desempeño superior al de los promedios de
	\textit{word embeddings} en tareas de clasificación de texto corto y mediano. Esa capacidad resulta
	crítica para distinguir los casos límite de la DDP, donde el acto jurídicamente relevante puede
	estar expresado en una sola oración dentro de un párrafo extenso.
	
	\subsection{Ingeniería de características híbridas}
	Aunque los \textit{embeddings} semánticos capturan el significado del texto, los modelos puramente
	textuales pueden aprender correlaciones léxicas espurias en lugar de la lógica jurídica subyacente.
	Por ello, una de las contribuciones centrales de este sistema es la arquitectura híbrida que combina
	la representación semántica con conocimiento normativo estructurado. Esa arquitectura opera en dos
	pasos: primero codifica los criterios del Artículo 4.\textsuperscript{o} como variables binarias y
	después las fusiona con el \textit{embedding} textual mediante una concatenación amplificada que
	equilibra ambas señales.
	
	\subsubsection{Codificación de variables normativas}
	Los cinco criterios de incompetencia del Artículo 4.\textsuperscript{o} se traducen en siete
	variables binarias ---los atributos jurídicos--- que codifican el análisis normativo de cada queja
	de forma estructurada. Esta codificación es el mecanismo que permite al clasificador razonar con
	base en la ley y no solo en la semántica del texto. Las variables y su descripción normativa se
	presentan en la Tabla~\ref{tab:atributos_juridicos}.
	
	\rowcolors{2}{rowblue}{white}
	
	\footnotesize
	\renewcommand{\arraystretch}{1.25}
	\setlength{\tabcolsep}{4pt}
	
	\begin{longtable}{|p{0.6cm}|p{5.1cm}|p{8.0cm}|}
		\caption{Atributos jurídicos estructurados derivados del Artículo 4.\textsuperscript{o}}
		\label{tab:atributos_juridicos} \\
		\hline
		\rowcolor{headerblue}
		\headercell{\#} & \headercell{Variable} & \headercell{Descripción normativa} \\
		\hline
		\endfirsthead
	
		\hline
		\rowcolor{headerblue}
		\headercell{\#} & \headercell{Variable} & \headercell{Descripción normativa} \\
		\hline
		\endhead
	
		1 & \texttt{es\_\allowbreak conflicto\_\allowbreak laboral\_\allowbreak art4\_\allowbreak I} &
		El fondo del conflicto es una relación laboral o sindical: salarios, plazas, escalafón o
		prestaciones. \\ \hline
	
		2 & \texttt{es\_\allowbreak evaluacion\_\allowbreak academica\_\allowbreak pura\_\allowbreak art4\_\allowbreak II} &
		La queja versa únicamente sobre la valoración del aprendizaje por parte de un profesor o de una
		comisión. \\ \hline
	
		3 & \texttt{hay\_\allowbreak vulneracion\_\allowbreak notoria\_\allowbreak en\_\allowbreak evaluacion} &
		Existe una violación de derechos dentro de un proceso evaluativo: extorsión, acoso o
		discriminación. \\ \hline
	
		4 & \texttt{es\_\allowbreak resolucion\_\allowbreak disciplinaria\_\allowbreak art4\_\allowbreak III} &
		La queja impugna una resolución disciplinaria ya emitida. \\ \hline
	
		5 & \texttt{es\_\allowbreak afectacion\_\allowbreak colectiva\_\allowbreak art4\_\allowbreak IV} &
		La afectación descrita es grupal y no individual. \\ \hline
	
		6 & \texttt{involucra\_\allowbreak oic\_\allowbreak art4\_\allowbreak V} &
		El asunto está siendo conocido por el Órgano Interno de Control. \\ \hline
	
		7 & \texttt{hay\_\allowbreak acto\_\allowbreak u\_\allowbreak omision\_\allowbreak autoridad} &
		Existe un acto u omisión de una autoridad institucional que vulnera derechos politécnicos. \\
		\hline
	\end{longtable}
	
	\normalsize
	\setlength{\tabcolsep}{6pt}
	\renewcommand{\arraystretch}{1}
	
	Estas variables se representan como un vector $\mathbf{f} \in \{0, 1\}^{7}$ y se determinan a partir
	del análisis del expediente, ya sea de forma manual durante la construcción del corpus o de forma
	automática, en una etapa futura, mediante un modelo extractor. La lógica de consistencia entre las
	variables y la clase asignada es la siguiente:
	
	\begin{itemize}
		\item Si la clase es No Competente por razón laboral,
		\texttt{es\_\allowbreak conflicto\_\allowbreak laboral\_\allowbreak art4\_\allowbreak I} $= 1$ y
		el resto de las variables $= 0$.
	
		\item Si la clase es Competente,
		\texttt{hay\_\allowbreak acto\_\allowbreak u\_\allowbreak omision\_\allowbreak autoridad} $= 1$ y
		todas las fracciones del Artículo 4.\textsuperscript{o} $= 0$.
	
		\item La variable 3 (\texttt{hay\_\allowbreak vulneracion\_\allowbreak notoria}) solo puede valer
		$1$ cuando la variable 2 (\texttt{es\_\allowbreak evaluacion\_\allowbreak academica\_\allowbreak pura})
		también vale $1$ y la clase es Competente; así se modela la excepción prevista en la fracción IV
		del Artículo 4.\textsuperscript{o}
	\end{itemize}
	
	\subsubsection{Concatenación y factor de amplificación ($\alpha$)}
	Codificadas las variables normativas, es necesario fusionarlas con el \textit{embedding} semántico
	de manera que ambas señales contribuyan de forma equitativa a la decisión del clasificador. La unión
	se realiza mediante concatenación horizontal junto con un factor de amplificación $\alpha$ que
	compensa la diferencia de escala entre los vectores:
	
	\begin{equation}
		\mathbf{x}_{\mathrm{final}} =
		\bigl[\, \mathbf{e} \,\big\|\, \alpha \cdot \mathbf{f} \,\bigr] \in \mathbb{R}^{d+7}
		\label{eq:concatenacion}
	\end{equation}
	
	donde $d \in \{300, 768\}$ según el vectorizador empleado.
	
	\paragraph{Justificación matemática del factor $\alpha$.} Sin amplificación ($\alpha = 1$), los
	siete atributos jurídicos quedan opacados por las $d$ dimensiones del \textit{embedding}, cuya norma
	promedio es bastante mayor. Para ilustrarlo, considérese la contribución relativa de cada componente
	a la norma del vector concatenado:
	
	\begin{equation}
		\frac{\lVert \alpha \cdot \mathbf{f} \rVert}
		     {\lVert \mathbf{e} \rVert + \lVert \alpha \cdot \mathbf{f} \rVert} =
		\frac{\alpha \sqrt{7}}{\lVert \mathbf{e} \rVert + \alpha \sqrt{7}}
		\label{eq:contribucion_alpha}
	\end{equation}
	
	Con $\lVert \mathbf{e} \rVert \approx 15$, valor típico para \textit{embeddings} de Sentence-BERT
	normalizados, y $\alpha = 1$:
	
	\[
		\frac{1 \cdot \sqrt{7}}{15 + \sqrt{7}} \approx \frac{2.65}{17.65} \approx 15\,\mathrm{\%}
	\]
	
	Con $\alpha = 3$:
	
	\[
		\frac{3 \sqrt{7}}{15 + 3\sqrt{7}} \approx \frac{7.94}{22.94} \approx 35\,\mathrm{\%}
	\]
	
	El valor $\alpha = 3.0$ se seleccionó de forma experimental como el punto de equilibrio en el que la
	señal jurídica pesa lo suficiente para influir en las fronteras de decisión del clasificador sin
	dominar por completo la información semántica del texto.
	
	\subsubsection{Simulación de producción: entrenamiento con metadatos e inferencia en modo autónomo}
	La arquitectura híbrida introduce una asimetría importante entre el escenario de entrenamiento y el
	de producción, que debe modelarse de forma explícita para evitar que el sistema resulte
	sobreoptimista en sus métricas de evaluación. Esa asimetría surge porque, durante la construcción
	del corpus, el personal jurídico anota manualmente los siete atributos booleanos del
	Artículo 4.\textsuperscript{o}, mientras que en producción el sistema recibe únicamente el relato
	libre del quejoso, sin anotación previa. La asimetría se modela mediante dos modos de operación:
	
	\begin{enumerate}
		\item \textbf{Modo asistido.} El personal de primer contacto completa el formulario con los
		atributos jurídicos mientras atiende al quejoso. El clasificador opera con el vector completo
		$\mathbf{x}_{\mathrm{final}} = [\mathbf{e} \,\|\, \alpha \cdot \mathbf{f}]$ y aprovecha tanto la
		semántica del texto como la señal normativa estructurada.
	
		\item \textbf{Modo autónomo.} Los atributos jurídicos se fijan en cero ($\mathbf{f} = \mathbf{0}$),
		de forma que el clasificador basa su veredicto exclusivamente en la representación semántica del
		texto: $\mathbf{x}_{\mathrm{final}} = [\mathbf{e} \,\|\, \mathbf{0}]$. Este modo representa el
		escenario real de producción y constituye la condición de estrés bajo la cual se evalúa la robustez
		de los modelos.
	\end{enumerate}
	
	La distinción entre ambos modos tiene implicaciones directas para la selección del clasificador
	óptimo: un modelo que dependa en exceso de la señal normativa amplificada puede exhibir métricas
	perfectas durante el entrenamiento y degradarse de forma significativa en producción, cuando esa
	señal no está disponible. Por esa razón, el criterio de selección del modelo final privilegia el
	desempeño en modo autónomo sobre el desempeño con metadatos.
	
	\subsection{Modelos de aprendizaje automático evaluados}
	Definido el vector de características híbridas, el siguiente paso es elegir el algoritmo que mejor
	aprenda a partir de él las fronteras de decisión legal. Para garantizar la selección del clasificador
	óptimo para el dominio jurídico de la DDP se evaluaron cuatro configuraciones de aprendizaje
	supervisado con paradigmas distintos de construcción de fronteras. Cada modelo aporta una
	perspectiva diferente sobre los datos, y su comparación permite identificar cuál se adapta mejor a
	la distribución de las quejas en el espacio vectorial.
	
	\subsubsection{Máquinas de soporte vectorial (SVM)}
	Las máquinas de soporte vectorial (\textit{Support Vector Machines}, SVM) son algoritmos de
	clasificación que buscan el hiperplano de máximo margen capaz de separar las clases en el espacio de
	características, lo que las hace especialmente robustas ante casos límite como las quejas situadas
	en la frontera entre competente e incompetente \cite{cortes1995svm}. Dado un conjunto de
	entrenamiento $\{(\mathbf{x}_i, y_i)\}_{i=1}^{n}$ con $y_i \in \{-1, +1\}$, la SVM resuelve el
	problema de optimización:
	
	\begin{equation}
		\begin{aligned}
			\min_{\mathbf{w}, b, \boldsymbol{\xi}} \quad
			& \frac{1}{2} \lVert \mathbf{w} \rVert^{2} + C \sum_{i=1}^{n} \xi_i \\
			\text{s.a.} \quad
			& y_i (\mathbf{w}^{\top} \mathbf{x}_i + b) \geq 1 - \xi_i, \qquad \xi_i \geq 0
		\end{aligned}
		\label{eq:svm}
	\end{equation}
	
	donde $C$ es el parámetro de regularización que controla la tolerancia al error de clasificación
	---con valor $C = 0.5$ en este sistema--- y $\boldsymbol{\xi}$ son las variables de holgura que
	permiten una clasificación blanda.
	
	\paragraph{Kernel RBF.} Cuando los datos no son linealmente separables en el espacio original, la
	SVM los proyecta a un espacio de mayor dimensión mediante una función de kernel. El kernel de base
	radial (\textit{Radial Basis Function}, RBF) que emplea este sistema calcula la similitud entre dos
	puntos como:
	
	\begin{equation}
		K(\mathbf{x}_i, \mathbf{x}_j) =
		\exp\!\bigl( -\gamma \lVert \mathbf{x}_i - \mathbf{x}_j \rVert^{2} \bigr)
		\label{eq:kernel_rbf}
	\end{equation}
	
	donde $\gamma$ controla el radio de influencia de cada vector de soporte. El kernel RBF resulta
	especialmente adecuado para \textit{embeddings} de alta dimensión porque puede construir fronteras
	de decisión no lineales arbitrariamente complejas, adaptándose a la distribución de los casos
	jurídicos en el espacio semántico \cite{cortes1995svm}.
	
	\paragraph{Ventaja en el dominio jurídico.} La SVM con kernel RBF maximiza el margen de separación
	entre las clases, lo que le confiere robustez natural ante casos límite: las quejas que caen en la
	frontera entre Competente y No Competente se clasifican según el hiperplano de máximo margen y no
	según la densidad local de puntos. Esa ventaja geométrica está condicionada, sin embargo, a que la
	señal de entrada sea suficientemente discriminativa; cuando los atributos jurídicos no están
	presentes, la eficacia del clasificador depende por completo de la calidad del \textit{embedding}
	semántico.
	
	\subsubsection{Regresión logística}
	A diferencia de la SVM, que construye una frontera geométrica de máximo margen, la regresión
	logística es un modelo lineal que estima directamente la probabilidad de pertenencia a cada clase.
	Ello le confiere una interpretabilidad valiosa en el contexto jurídico, donde la explicabilidad de
	la decisión puede exigirse institucionalmente. La probabilidad de pertenencia a la clase positiva se
	estima mediante la función logística o sigmoide:
	
	\begin{equation}
		P(y = 1 \mid \mathbf{x}) = \sigma(\mathbf{w}^{\top} \mathbf{x} + b) =
		\frac{1}{1 + e^{-(\mathbf{w}^{\top} \mathbf{x} + b)}}
		\label{eq:sigmoide}
	\end{equation}
	
	Los parámetros $\mathbf{w}$ y $b$ se estiman minimizando la función de pérdida de entropía cruzada
	binaria:
	
	\begin{equation}
		\mathcal{L}(\mathbf{w}, b) =
		-\frac{1}{n} \sum_{i=1}^{n}
		\Bigl[\, y_i \log \hat{p}_i + (1 - y_i) \log (1 - \hat{p}_i) \,\Bigr]
		+ \lambda \lVert \mathbf{w} \rVert^{2}
		\label{eq:entropia_cruzada}
	\end{equation}
	
	donde $\lambda$ es el término de regularización L2. La configuración empleada en este sistema es
	\texttt{max\_iter=1000} con $C = 0.3$. Los coeficientes $w_j$ son además directamente
	interpretables: un coeficiente positivo asociado a una dimensión del \textit{embedding} indica que
	ese componente favorece la clasificación como Competente, lo que facilita la auditoría del modelo.
	
	\subsubsection{Ensamble de árboles: Random Forest}
	Donde la SVM y la regresión logística construyen una sola frontera de decisión, Random Forest
	resuelve el problema agregando múltiples árboles de decisión entrenados sobre subconjuntos
	aleatorios de los datos. Esta estrategia de ensamble reduce la varianza de la predicción individual
	y confiere al modelo una notable capacidad de generalización \cite{breiman2001rf}. Formalmente:
	
	\begin{equation}
		\hat{y}_{\mathrm{RF}}(\mathbf{x}) = \frac{1}{B} \sum_{b=1}^{B} T_b(\mathbf{x})
		\label{eq:random_forest}
	\end{equation}
	
	donde $B$ es el número de árboles ---$B = 300$ en este sistema--- y $T_b$ es el $b$-ésimo árbol
	entrenado sobre una muestra \textit{bootstrap}. Los árboles particionan el espacio de
	características con cortes ortogonales a los ejes y construyen de forma recursiva regiones de
	decisión rectangulares. Aunque en espacios de dimensión muy alta esa estrategia puede ser menos
	eficiente que el hiperplano de máximo margen de la SVM, el promedio de un número grande de árboles
	con submuestreo aleatorio de características reduce de forma significativa la varianza y resulta
	especialmente robusto cuando la representación semántica del texto no está apoyada por atributos
	estructurados.
	
	\subsubsection{Ensamble por votación suave}
	Dado que cada clasificador individual tiene fortalezas complementarias, una estrategia de ensamble
	que combine sus predicciones puede alcanzar un desempeño superior al de cualquiera por separado. La
	votación suave (\textit{soft voting}) aprovecha esa propiedad promediando las probabilidades
	predichas por múltiples clasificadores base para obtener una predicción final más estable
	\cite{dietterich2000ensemble}:
	
	\begin{equation}
		\hat{y} = \arg\max_{c} \frac{1}{M} \sum_{m=1}^{M} p_m(c \mid \mathbf{x})
		\label{eq:votacion_suave}
	\end{equation}
	
	donde $M$ es el número de modelos base y $p_m(c \mid \mathbf{x})$ es la probabilidad que el modelo
	$m$ asigna a la clase $c$. A diferencia de la votación dura, que solo cuenta votos de clase, la
	votación suave aprovecha la magnitud de la confianza de cada clasificador, lo que suele producir
	predicciones mejor calibradas cuando las tasas de error de los modelos base no están
	correlacionadas.
	
	En este sistema el ensamble integra SVM-RBF, regresión logística y Random Forest. La motivación es
	que cada uno aporta algo distinto: la SVM maximiza el margen geométrico, la regresión logística
	produce probabilidades linealmente calibradas y Random Forest captura fronteras no lineales por
	partición del espacio. Si sus errores no están correlacionados, el ensamble reduce la varianza de la
	predicción respecto a cualquiera de los modelos base. Para que \texttt{VotingClassifier} opere en
	modo \texttt{voting="soft"}, todos los estimadores deben exponer el método
	\texttt{predict\_proba}; dado que \texttt{SVC} no lo implementa de forma nativa, se envuelve en
	\texttt{CalibratedClassifierCV}, que estima probabilidades mediante validación cruzada interna sobre
	el conjunto de entrenamiento \cite{pedregosa2011sklearn}.
	
	\subsection{Estrategias de entrenamiento y validación robusta}
	\label{sec:validacion_robusta}
	La selección del mejor clasificador no puede basarse en métricas obtenidas sobre los mismos datos
	con que se entrenó, pues ello produciría una estimación optimista del desempeño real. Las
	estrategias de entrenamiento y validación robusta garantizan que las métricas reportadas reflejen la
	capacidad de generalización del sistema y no sean un artefacto de una partición particular de los
	datos.
	
	\subsubsection{Validación cruzada $K$-fold estratificada}
	La validación cruzada $K$-fold es la estrategia central para estimar el desempeño de los modelos de
	manera confiable: divide el conjunto de datos en $K$ particiones de tamaño aproximadamente igual y
	rota el conjunto de evaluación, de modo que todos los datos se usen para prueba exactamente una vez.
	En cada iteración, $K - 1$ particiones sirven para entrenamiento y la restante para evaluación, lo
	que produce una estimación promedio del desempeño:
	
	\begin{equation}
		F1_{\mathrm{CV}} = \frac{1}{K} \sum_{k=1}^{K} F1_k
		\label{eq:f1_cv}
	\end{equation}
	
	La variante estratificada garantiza que la proporción de clases se mantenga en cada pliegue,
	condición crítica ante corpus con posible desbalance entre quejas competentes e incompetentes
	\cite{pedregosa2011sklearn}. En este proyecto se empleó $K = 5$, lo que significa que en cada
	iteración el 80\,\% de los datos se usa para entrenamiento y el 20\,\% restante para evaluación.
	Además del promedio, la desviación estándar entre pliegues ($\sigma_{\mathrm{CV}}$) se emplea como
	medida de estabilidad del modelo:
	
	\begin{equation}
		\sigma_{\mathrm{CV}} =
		\sqrt{\frac{1}{K} \sum_{k=1}^{K} \bigl( F1_k - F1_{\mathrm{CV}} \bigr)^{2}}
		\label{eq:sigma_cv}
	\end{equation}
	
	Un $\sigma_{\mathrm{CV}}$ cercano a cero indica que el modelo no es sensible a la partición
	específica de los datos, condición deseable para un sistema de apoyo en producción.
	
	\paragraph{Nota de implementación.} En este proyecto la validación cruzada con $K = 5$ se aplica
	exclusivamente sobre el conjunto de entrenamiento ---el 70\,\% del corpus, es decir, 367 casos--- y
	no sobre el conjunto de datos completo, de modo que cada pliegue evalúa alrededor de 73 casos
	internos del entrenamiento. Esta validación se complementa con un conjunto \textit{holdout}
	independiente del 15\,\% (79 casos), aislado por completo durante el entrenamiento y la selección de
	hiperparámetros, cuyo F1 ponderado ---evaluado en modo autónomo, sin atributos jurídicos---
	constituye el criterio definitivo de selección del modelo.
	
	\subsubsection{Evaluación mediante métricas de desempeño}
	Las métricas de evaluación deben reflejar las prioridades institucionales de la DDP y no únicamente
	la exactitud global del modelo. Se emplean, por tanto, métricas que distinguen entre los distintos
	tipos de error y asignan mayor peso al de mayor costo institucional: clasificar como incompetente
	una queja que sí debía atenderse. Para comprenderlas conviene establecer primero los cuatro
	resultados posibles de una clasificación binaria, organizados en la matriz de confusión:
	
	\begin{itemize}
		\item \textbf{Verdadero positivo (VP):} la queja es competente y el modelo la clasifica como
		competente.
		\item \textbf{Verdadero negativo (VN):} la queja es incompetente y el modelo la clasifica como
		incompetente.
		\item \textbf{Falso positivo (FP):} la queja es incompetente pero el modelo la clasifica como
		competente.
		\item \textbf{Falso negativo (FN):} la queja es competente pero el modelo la clasifica como
		incompetente. En el contexto de la DDP este es el error más grave, pues implica descartar una queja
		que debía ser atendida.
	\end{itemize}
	
	\paragraph{Exactitud.} La exactitud (\textit{accuracy}) mide la proporción de predicciones correctas
	sobre el total de instancias:
	
	\begin{equation}
		\text{Exactitud} = \frac{VP + VN}{VP + VN + FP + FN}
		\label{eq:accuracy}
	\end{equation}
	
	En el contexto de la DDP, donde las consecuencias de un falso negativo son más graves que las de un
	falso positivo, la exactitud por sí sola no basta. Por ello se complementa con las métricas de
	precisión, exhaustividad y F1.
	
	\paragraph{Precisión, exhaustividad y F1.} La precisión mide qué proporción de las quejas
	clasificadas como competentes lo son realmente y penaliza las falsas alarmas:
	
	\begin{equation}
		\text{Precisión} = \frac{VP}{VP + FP}
		\label{eq:precision}
	\end{equation}
	
	La exhaustividad (\textit{recall}) mide qué proporción de las quejas realmente competentes fue
	identificada de forma correcta y penaliza los casos omitidos:
	
	\begin{equation}
		\text{Exhaustividad} = \frac{VP}{VP + FN}
		\label{eq:recall}
	\end{equation}
	
	En el contexto de la DDP esta métrica tiene mayor peso institucional: una exhaustividad baja implica
	que quejas legítimas se están descartando sin revisión. El F1 es la media armónica entre ambas y
	ofrece un balance que penaliza los extremos donde una métrica es alta a costa de la otra:
	
	\begin{equation}
		F1 = \frac{2 \cdot \text{Precisión} \cdot \text{Exhaustividad}}
		          {\text{Precisión} + \text{Exhaustividad}}
		\label{eq:f1}
	\end{equation}
	
	El F1 ponderado (\textit{weighted F1}) extiende el cálculo a ambas clases, ponderando por el número
	de instancias de cada una:
	
	\begin{equation}
		F1_{w} = \sum_{c \in C} \frac{n_c}{N} \cdot F1_c
		\label{eq:f1_ponderado}
	\end{equation}
	
	donde $n_c$ es el número de instancias de la clase $c$ y $N$ el total. Esta métrica es la principal
	para la selección del modelo óptimo en este sistema.
	
	\paragraph{Análisis de la varianza.} La desviación estándar del F1 en validación cruzada
	($\sigma_{\mathrm{CV}}$) complementa el valor medio al revelar la estabilidad del modelo. Un modelo
	con $F1_{\mathrm{CV}} = 0.98$ y $\sigma_{\mathrm{CV}} = 0.04$ puede resultar menos confiable que uno
	con $F1_{\mathrm{CV}} = 0.97$ y $\sigma_{\mathrm{CV}} = 0.00$, pues el primero es sensible a la
	distribución particular de los datos en cada pliegue.
	
	\subsubsection{Carga perezosa y serialización con joblib}
	Seleccionado el modelo óptimo, este debe almacenarse y ponerse en producción de forma que garantice
	la reproducibilidad de sus predicciones y minimice el tiempo de arranque del sistema. Para ello se
	combinan la carga perezosa y la serialización con \texttt{joblib}, dos patrones de diseño que
	optimizan la eficiencia operativa del clasificador.
	
	La carga perezosa (\textit{lazy loading}) inicializa los recursos computacionalmente costosos ---como
	el modelo de Sentence-BERT o el de spaCy--- únicamente cuando se requieren por primera vez, y no al
	inicio del programa. La serialización, por su parte, convierte el objeto de Python completo
	---vectorizador y clasificador juntos--- en un archivo binario \texttt{.pkl}:
	
	\begin{center}
		\small
		\texttt{Pipeline(vectorizador, clasificador)}
		$\xrightarrow{\ \texttt{joblib.dump}\ }$
		\texttt{modelo.pkl}
		$\xrightarrow{\ \texttt{joblib.load}\ }$
		Pipeline lista para inferencia
	\end{center}
	
	Esta estrategia es técnicamente relevante por tres razones:
	
	\begin{enumerate}
		\item \textbf{Integridad.} Vectorizador y clasificador se serializan juntos, lo que elimina el
		riesgo de usar un clasificador entrenado con un vectorizador distinto al de producción.
	
		\item \textbf{Eficiencia.} \texttt{joblib} emplea compresión paralela y mapeo de memoria
		(\textit{memory mapping}), por lo que resulta bastante más rápido que \texttt{pickle} estándar para
		arreglos numéricos de gran tamaño, como los parámetros internos de Sentence-BERT.
	
		\item \textbf{Reproducibilidad.} El archivo \texttt{.pkl} incluye la semilla aleatoria
		(\texttt{random\_state=42}) y todos los hiperparámetros, lo que garantiza predicciones idénticas
		entre sesiones y entornos de ejecución.
	\end{enumerate}
	
	La convención de nomenclatura de los archivos serializados es:
	
	\begin{center}
		\texttt{<vectorizador>\_\_<clasificador>\_\_<timestamp>.pkl}
	\end{center}
	
	lo que permite la trazabilidad temporal de los modelos entrenados y facilita la comparación y
	auditoría de versiones en el contexto institucional de la DDP.
